% Source of the listings (read from the repository at compile time, not copied):
%   App Store collector   scripts/appstore-scraper/v006/    scraper.py, input.json
%   Google Play collector scripts/googleplay-scraper/v002/  scraper.py, input.json, requirements.txt
% To show a newer version, change the paths below and the two comments above.
\chapter{Scraping Scripts}
\label{app:scripts}
Both collectors read a JSON file and print the result as JSON on standard
output. They record what each store returns for a keyword search or, for Google
Play, for a list of package IDs. They apply no eligibility rule and no ranking,
and they do not observe app behaviour. The listings are the files used in the
project, shown unchanged.

\section{Input Format}
\begin{description}
  \item[\texttt{keywords}] Search terms, one request each (both scripts).
  \item[\texttt{ids}] Google Play package IDs to look up directly, without a search.
  \item[\texttt{region}] Two-letter storefront; \texttt{de} selects the German store.
  \item[\texttt{language}] Interface language requested from Google Play.
  \item[\texttt{limit}] Maximum results per keyword (default 20 on Google Play).
\end{description}

\lstinputlisting[style=jsoncode, caption={Example input for the App Store collector.}, label={lst:input-apple}]{../scripts/appstore-scraper/v006/input.json}

\lstinputlisting[style=jsoncode, caption={Example input for the Google Play collector.}, label={lst:input-play}]{../scripts/googleplay-scraper/v002/input.json}

\section{App Store Collector}
The script calls Apple's documented iTunes Search API once per keyword and
storefront, keeps iPhone-compatible results in the order returned, and pauses
between requests to respect the rate limit. Rating counts and averages describe
the requested storefront (Table~\ref{tab:storefront-fields}).

\lstinputlisting[style=pythoncode, caption={App Store collector.}, label={lst:scraper-apple}]{../scripts/appstore-scraper/v006/scraper.py}

\section{Google Play Collector}
The script uses the open-source \texttt{google-play-scraper} parser (version
1.2.7) to read the embedded data of the search page and the product page of each
app. It restores TLS certificate verification, which the library disables on
import, and rejects any request or redirect that drops the storefront
parameter, so a result cannot silently come from another country. Failed
lookups are recorded under \texttt{errors}. Search positions follow the order of
the search response and are not verified as the store's native ranking. Rating
counts, rating distribution and install figures are worldwide on Google Play
(Table~\ref{tab:storefront-fields}).

\lstinputlisting[style=pythoncode, caption={Google Play collector.}, label={lst:scraper-play}]{../scripts/googleplay-scraper/v002/scraper.py}
